\documentclass[10pt,a4paper]{amsart}
\usepackage{amsmath,amssymb,amsthm,graphicx}
\usepackage{verbatim}
\newenvironment{reference}{\comment}{\endcomment}
%\newenvironment{reference}{}{}
\newenvironment{slogan}{\comment}{\endcomment}
\newenvironment{history}{\comment}{\endcomment}

% For commutative diagrams we use Xy-pic
\usepackage[all]{xy}

% We use 2cell for 2-commutative diagrams.
\xyoption{2cell}
% Theorem environments.
%
\theoremstyle{plain}
\newtheorem{theorem}[subsection]{Theorem}
\newtheorem{proposition}[subsection]{Proposition}
\newtheorem{lemma}[subsection]{Lemma}

\theoremstyle{definition}
\newtheorem{definition}[subsection]{Definition}
\newtheorem{example}[subsection]{Example}
\newtheorem{exercise}[subsection]{Exercise}
\newtheorem{situation}[subsection]{Situation}

\theoremstyle{remark}
\newtheorem{remark}[subsection]{Remark}
\newtheorem{remarks}[subsection]{Remarks}

\numberwithin{equation}{subsection}

% Macros
%
\def\lim{\mathop{\mathrm{lim}}\nolimits}
\def\colim{\mathop{\mathrm{colim}}\nolimits}
\def\Spec{\mathop{\mathrm{Spec}}}
\def\Hom{\mathop{\mathrm{Hom}}\nolimits}
\def\Ext{\mathop{\mathrm{Ext}}\nolimits}
\def\SheafHom{\mathop{\mathcal{H}\!\mathit{om}}\nolimits}
\def\SheafExt{\mathop{\mathcal{E}\!\mathit{xt}}\nolimits}
\def\Sch{\mathit{Sch}}
\def\Mor{\mathop{\mathrm{Mor}}\nolimits}
\def\Ob{\mathop{\mathrm{Ob}}\nolimits}
\def\Sh{\mathop{\mathit{Sh}}\nolimits}
\def\NL{\mathop{N\!L}\nolimits}
\def\CH{\mathop{\mathrm{CH}}\nolimits}
\def\proetale{{pro\text{-}\acute{e}tale}}
\def\etale{{\acute{e}tale}}
\def\QCoh{\mathit{QCoh}}
\def\Ker{\mathop{\mathrm{Ker}}}
\def\Im{\mathop{\mathrm{Im}}}
\def\Coker{\mathop{\mathrm{Coker}}}
\def\Coim{\mathop{\mathrm{Coim}}}

% Boxtimes
%
\DeclareMathSymbol{\boxtimes}{\mathbin}{AMSa}{"02}

%
% Macros for moduli stacks/spaces
%
\def\QCohstack{\mathcal{QC}\!\mathit{oh}}
\def\Cohstack{\mathcal{C}\!\mathit{oh}}
\def\Spacesstack{\mathcal{S}\!\mathit{paces}}
\def\Quotfunctor{\mathrm{Quot}}
\def\Hilbfunctor{\mathrm{Hilb}}
\def\Curvesstack{\mathcal{C}\!\mathit{urves}}
\def\Polarizedstack{\mathcal{P}\!\mathit{olarized}}
\def\Complexesstack{\mathcal{C}\!\mathit{omplexes}}
% \Pic is the operator that assigns to X its picard group, usage \Pic(X)
% \Picardstack_{X/B} denotes the Picard stack of X over B
% \Picardfunctor_{X/B} denotes the Picard functor of X over B
\def\Pic{\mathop{\mathrm{Pic}}\nolimits}
\def\Picardstack{\mathcal{P}\!\mathit{ic}}
\def\Picardfunctor{\mathrm{Pic}}
\def\Deformationcategory{\mathcal{D}\!\mathit{ef}}
\title{A termwise split sequence of complexes of complexes gives a connecting morphism}
\author{The Stacks Project Authors}
\date{stacks-project:homology.tex:local-snapshot-20261012; GNU FDL 1.2 or later}
\begin{document}
\maketitle
\noindent\textit{Editorial note.} All delivered mathematical text is the author's own, copied from the pinned Stacks Project chapter file homology.tex; no statement or proof step was rewritten, reordered or removed.\par\medskip
\noindent\textit{Editorial note.} The proof is the author's complete argument; both identities it uses are displayed inside the counted span and the closing words 'as desired' follow the computation that establishes the claim, so no step is deferred.\par\medskip
\noindent\textit{Editorial note.} Every labelled reference occurring inside the counted spans is carried in assembly.external\_references as a non-load-bearing navigation pointer (same-chapter ones at their definition site, cross-chapter ones at the invocation inside the counted span), so no load-bearing mathematics is replaced by a stub. No bibliography citation occurs inside the counted spans.\par\medskip
\noindent\textit{Editorial note.} The same-name check for this card scanned every staged Stacks chapter of this repository -- all sources/stacks-*/pinned-original-source.tex files, the card's own chapter excluded -- and the labels found are recorded in dedup.same\_name\_cross\_chapter on the card.\par\medskip
\noindent\textit{Editorial note.} The preamble is the chapter's own native preamble excerpt already used for this source: the comment-type reference, slogan and history environments (native preamble lines 8-13), the Xy-pic setup (native preamble lines 14-18) and the native theorem declarations with the mathematical macros (native preamble lines 154-217, the actual last line of the file).\par\medskip
\noindent\textit{Editorial note.} Same-name and twin check: the chapter also states lemma-ses-termwise-split at source lines 3578-3598, which constructs the same connecting morphism in the homological grading (a morphism C to A[-1]). The two statements have the same conclusion under the reindexing A\_n = A\textasciicircum{}\{-n\} and equivalent splitting hypotheses, so per the lane rule for twins only one of them is counted, namely this cohomological version. The homological twin is recorded in lane-s4/DISPOSITIONS-stacks-homology.json and is not counted anywhere.\par\medskip
\noindent\textit{Allowed standard prerequisites (editorial).}\begin{itemize}
\item Splittings of a complex of complexes. Same chapter, the discussion of a complex of complexes and of compatible direct sum splittings of its terms. Supplies the maps s\textasciicircum{}n and pi\textasciicircum{}n induced by the splitting B\textasciicircum{}n = A\textasciicircum{}n plus C\textasciicircum{}n that the construction uses.
\item Differentials and the shift of a complex. Same chapter, the definition of the differential of a complex and of the shift A[1]. Supplies the objects C\textasciicircum{}n and A[1]\textasciicircum{}\{n+1\} connected by the constructed family and the sign convention used in the computation.
\end{itemize}\medskip
\begin{lemma}
\label{lemma-ses-termwise-split-cochain}
Let $\mathcal{A}$ be an additive category.
Let
$$
0 \to A^\bullet \to B^\bullet \to C^\bullet \to 0
$$
be a complex (!) of complexes.
Suppose that we are given splittings $B^n = A^n \oplus C^n$
compatible with the maps in the displayed sequence.
Let $s^n : C^n \to B^n$ and $\pi^n : B^n \to A^n$ be the
corresponding maps. Then the family of morphisms
$$
\pi^{n + 1} \circ d_B^n \circ s^n
:
C^n \to A^{n + 1}
$$
define a morphism of complexes $\delta : C^\bullet \to A[1]^\bullet$.
\end{lemma}

\begin{proof}
Denote $i : A^\bullet \to B^\bullet$ and $q : B^\bullet \to C^\bullet$
the maps of complexes in the short exact sequence. Then
$i^{n + 1} \circ \pi^{n + 1} \circ d_B^n \circ s^n =
d_B^n \circ s^n - s^{n + 1} \circ d_C^n$. Hence
$i^{n + 2} \circ d_A^{n + 1} \circ \pi^{n + 1} \circ d_B^n \circ s^n =
d_B^{n + 1} \circ (d_B^n \circ s^n - s^{n + 1} \circ d_C^n) =
- d_B^{n + 1} \circ s^{n + 1} \circ d_C^n$ as desired.
\end{proof}

\end{document}
